\documentclass[12pt]{article}
\usepackage[utf8]{inputenc}
\usepackage[T1]{fontenc}
\usepackage{hyperref}
\title{Nat\"ugu (Retired) Language Resources for K-12 Curriculum Development}
\author{Compiled by Claude Sonnet 5.5 (claude-sonnet-5-5) for Andr\'as Kornai}
\date{}
\begin{document}
\maketitle

\section{Introduction}
The code stc (``Natügu (Retired)'' in Glottolog's bookkeeping family) is a retired ISO 639-3 code. According to the search results, stc (Santa Cruz) was retired in 2009 and split into Natügu [ntu] and Nal\"ogo [nlz]. Current resources should therefore be sought under ntu. Natügu is the main language of Nend\"o (Santa Cruz Island) in the Solomon Islands. Sources give about 5,000 to 5,900 speakers. It is Oceanic Austronesian, though it was once thought to be Papuan. Only general reference pages were found in the search.

\section{Primary Educational Resources}
No primers, textbooks or school materials were found in the search.

\section{Linguistic and Reference Materials}
\begin{itemize}
\item Wikipedia article on the Natügu language: \url{https://en.wikipedia.org/wiki/Nat%C3%BCgu_language}
\item Wikipedia article on the Santa Cruz language: \url{https://en.wikipedia.org/wiki/Santa_Cruz_language}
\item Glottolog entry for Natügu (Retired): \url{https://glottolog.org/resource/languoid/id/natu1249}
\item Omniglot page on Santa Cruz: \url{https://omniglot.com/writing/santacruz.htm}
\item WALS entry (code nbk, which I take to be a related Santa Cruz-area language; not checked): \url{https://wals.info/languoid/lect/wals_code_nbk}
\end{itemize}

\section{Religious and Community Resources}
None found.

\section{Digital and Audio Resources}
\begin{itemize}
\item PARADISEC catalogue item that appeared in the results (contents not checked): \url{https://catalog.paradisec.org.au/repository/SAW2/077}
\end{itemize}

\section{Curriculum Development Context}
Because stc is retired, curriculum planning should target Natügu (ntu) and the neighbouring Nal\"ogo (nlz), and should not be built on stc. The search found no school materials. The PARADISEC catalogue and the reference pages above are starting points for tracing documentation and speakers.

\end{document}